\documentclass[11pt,letterpaper]{article}
\usepackage[margin=0.7in]{geometry}
\usepackage[T1]{fontenc}
\usepackage{lmodern}
\usepackage{amsmath,amssymb,booktabs,array,xcolor,listings,hyperref,fancyhdr}
\definecolor{accent}{RGB}{25,65,100}
\hypersetup{colorlinks=true,linkcolor=accent,urlcolor=accent}
\setlength{\parindent}{0pt}
\setlength{\parskip}{4pt}
\setlength{\headheight}{14pt}
\pagestyle{fancy}\fancyhf{}\fancyhead[L]{\small CSI2110 -- Study Notes}\fancyhead[R]{\small Fall 2026}\fancyfoot[C]{\thepage}
\lstset{basicstyle=\small\ttfamily,breaklines=true,columns=fullflexible,keepspaces=true,frame=single,rulecolor=\color{accent!25},backgroundcolor=\color{accent!3},showstringspaces=false,tabsize=2}
\newcommand{\sourcepages}[1]{\textit{Source: PDF pages #1.}}
\newcommand{\keyidea}[1]{\par\textbf{Key idea.} #1\par}
\newcommand{\cost}[1]{\ensuremath{O(#1)}}
\newcommand{\method}[1]{\texttt{#1}}

\title{\textbf{Lecture 4: Trees and Binary Trees}\\\large CSI2110 -- Data Structures and Algorithms}\author{Lecture study notes}\date{}\begin{document}\maketitle
\textbf{Source document:} \path{../lec04-tree and binary tree.pdf}. These notes synthesize the source and add worked explanations. References use PDF page numbers. Pseudocode includes boundary cases.
\tableofcontents\newpage

\section{Graphs, trees, and rooted terminology}
\sourcepages{5--11}
A graph $G=(V,E)$ has vertices and edges. An undirected edge joins two vertices. A path follows adjacent vertices; a simple path does not repeat vertices. A cycle closes a path without otherwise repeating vertices. A connected graph has a path between every pair of vertices.

A \textbf{tree} is a connected, acyclic undirected graph. Selecting a root determines parent/child relationships. The root has no parent; every other node has exactly one parent.
\begin{itemize}
\item Siblings share a parent; ancestors are parents, grandparents, etc.; descendants are children, grandchildren, etc.
\item An internal node has at least one child; an external node (leaf) has none.
\item The subtree rooted at $v$ contains $v$ and all its descendants.
\item Distance counts edges. The root has depth zero. A node's depth is its distance from the root.
\item A nonempty tree's height is its maximum node depth, or its longest root-to-leaf path length in edges. A one-node tree has height zero.
\end{itemize}

\section{Binary-tree variants}
\sourcepages{13--17}
\begin{center}\begin{tabular}{@{}lp{4.8in}@{}}\toprule
Type & Condition\\\midrule
Binary & At most two children per node, distinguished as left and right.\\
Full (proper) & Every node has either zero or two children.\\
Perfect & Full, with all leaves at the same depth.\\
Complete & All levels before the last are full; the last is filled left to right.\\\bottomrule
\end{tabular}\end{center}
A perfect tree is both full and complete. A full tree need not be complete or balanced. A complete tree may have one node with only a left child. Decision trees and expression trees with binary operators are natural examples of full trees.

\section{Counting and height properties}
\sourcepages{19--33}
Let $n$ be total nodes, $n_I$ internal nodes, $n_E$ leaves, and $h$ height. Assume nonempty trees and base-two logarithms.

\subsection{Any tree}
There is exactly one \emph{simple} path between every pair of nodes. Connectivity gives existence; two distinct simple paths would imply a cycle. A tree has $n-1$ edges because every non-root node contributes one parent edge.

\subsection{Any binary tree}
At level $i$ there are at most $2^i$ nodes. Therefore,
\begin{align*}
h+1&\leq n\leq2^{h+1}-1,\\
1&\leq n_E\leq2^h,\\
h&\leq n_I\leq2^h-1,\\
\lceil\log_2(n+1)\rceil-1&\leq h\leq n-1.
\end{align*}
A chain attains $n=h+1$; a perfect tree attains $n=2^{h+1}-1$. The ceiling accounts for integer height and strengthens the real lower bound $\log_2(n+1)-1$.

\subsection{Full binary tree}
Each internal node contributes two child edges:
\[
2n_I=n-1=n_I+n_E-1\quad\Longrightarrow\quad n_E=n_I+1.
\]
Thus $n$ is odd, $n_E=(n+1)/2$, and $n_I=(n-1)/2$.
\begin{align*}
2h+1&\leq n\leq2^{h+1}-1,\\
h+1&\leq n_E\leq2^h,\\
h&\leq n_I\leq2^h-1,\\
\lceil\log_2(n+1)\rceil-1&\leq h\leq(n-1)/2.
\end{align*}
A full tree can still have height $\Theta(n)$: a long internal spine with a leaf sibling at each internal node.

\subsection{Perfect and complete trees}
For a perfect tree,
\[
n=2^{h+1}-1,\quad n_E=2^h,\quad n_I=2^h-1,\quad
h=\log_2(n+1)-1.
\]
For a complete tree,
\[
2^h\leq n\leq2^{h+1}-1,\qquad h=\lfloor\log_2n\rfloor.
\]
The height formula follows from $h\leq\log_2n<h+1$. From its array layout, a complete tree also has $n_I=\lfloor n/2\rfloor$ and $n_E=\lceil n/2\rceil$.

\textbf{Example:} with $h=3$, a binary tree has $4$--$15$ nodes; a full tree $7$--$15$; a complete tree $8$--$15$; a perfect tree exactly $15$ nodes and $8$ leaves.

\section{Tree ADTs}
\sourcepages{35--36}
General tree accessors include root, parent, and children. Queries test whether a position is root, internal, or external. Other operations report size, iterate elements/positions, replace an element, or swap elements. Updates depend on the application.

A binary-tree ADT adds left-child, right-child, and sibling accessors and specialized updates such as expanding a leaf or removing a leaf together with its parent.

\section{Traversal orders}
\sourcepages{37--59,65--68}
A traversal visits every node in a specified order; process ordered children left to right.
\begin{center}\begin{tabular}{@{}ll@{}}\toprule
Traversal & Binary-tree order\\\midrule
Breadth-first & Level by level, using a queue\\
Preorder & Root, left subtree, right subtree\\
Inorder & Left subtree, root, right subtree\\
Postorder & Left subtree, right subtree, root\\\bottomrule
\end{tabular}\end{center}
Preorder and postorder apply to arbitrary rooted trees; the lecture's inorder is for binary trees.

\subsection{Breadth-first traversal}
\par\noindent\begin{minipage}{\linewidth}
\begin{lstlisting}
BFS(root):
    if root == null: return
    enqueue(root)
    while queue is not empty:
        v = dequeue()
        visit(v)
        for each child w of v, in order:
            enqueue(w)
\end{lstlisting}
\end{minipage}\par
Each node is enqueued and dequeued once. Time is $\Theta(n)$ for constant-time visits; auxiliary space is $O(w)$ for maximum level width $w$, hence $O(n)$ worst case and $O(1)$ for a chain.

\subsection{Recursive depth-first traversal}
\par\noindent\begin{minipage}{\linewidth}
\begin{lstlisting}
preorder(v):
    if v == null: return
    visit(v)
    preorder(v.left)
    preorder(v.right)

inorder(v):
    if v == null: return
    inorder(v.left)
    visit(v)
    inorder(v.right)

postorder(v):
    if v == null: return
    postorder(v.left)
    postorder(v.right)
    visit(v)
\end{lstlisting}
\end{minipage}\par
Null checks correctly handle missing children. Time is $\Theta(n)$; auxiliary recursion space is $O(h+1)$, which is $O(\log n)$ for a complete tree and $O(n)$ for a chain. Output storage is separate.

\textbf{Lecture example (PDF page 46):} A has children B and G; B has children C and D; D has children E and F; G has children H and I.
\begin{itemize}
\item BFS: A, B, G, C, D, H, I, E, F.
\item Preorder: A, B, C, D, E, F, G, H, I.
\item Inorder: C, B, E, D, F, A, H, G, I.
\item Postorder: C, E, F, D, B, H, I, G, A.
\end{itemize}

\subsection{Iterative preorder and Euler tour}
\par\noindent\begin{minipage}{\linewidth}
\begin{lstlisting}
if root != null: push(root)
while stack is not empty:
    v = pop()
    visit(v)
    if v.right != null: push(v.right)
    if v.left != null: push(v.left)
\end{lstlisting}
\end{minipage}\par
Push the right child first so the left child is popped first. The pending-node stack uses $O(h+1)$ space in a binary tree.

An Euler tour offers three action points: before traversing the left subtree (preorder), between left and right (inorder), and after the right subtree (postorder). This unifies the three depth-first orders.

\section{Traversal applications}
\sourcepages{42--43,60--64}
\subsection{Structured documents and directory sizes}
Preorder prints headings before subheadings, with depth determining indentation. Postorder combines child sizes before calculating a directory total. The slide's files $10,25,20,3,2,1$ KB total $61$ KB, ignoring directory overhead.

\subsection{Subtree height}
Let a missing subtree have height $-1$, giving leaves height zero:
\par\noindent\begin{minipage}{\linewidth}
\begin{lstlisting}
height(v):
    if v == null: return -1
    return 1 + max(height(v.left), height(v.right))
\end{lstlisting}
\end{minipage}\par
This postorder algorithm takes $\Theta(n)$ time for the queried subtree and $O(h+1)$ auxiliary space.

\subsection{Expression trees}
Leaves are operands; internal nodes are binary operators. Preorder gives prefix notation, postorder postfix. Plain inorder can lose grouping; print parentheses around internal subtrees to preserve the expression.
\par\noindent\begin{minipage}{\linewidth}
\begin{lstlisting}
printExpr(v):
    if v is a leaf: print(v.element); return
    print("(")
    printExpr(v.left)
    print(v.operator)
    printExpr(v.right)
    print(")")

evalExpr(v):
    if v is a leaf: return its numeric value
    x = evalExpr(v.left)
    y = evalExpr(v.right)
    return apply(v.operator, x, y)
\end{lstlisting}
\end{minipage}\par
For $((2\times(a-1))+(3\times b))$, postfix is \method{2 a 1 - * 3 b * +}. With $a=5,b=2$, the value is $14$. Time is $\Theta(n)$ with constant-time operator evaluation or output per token.

\section{Representations and updates}
\sourcepages{70--81}
\subsection{Linked binary tree}
Each node stores an element, parent, left child, and right child. Given valid nodes, local access, replacement, and element swapping cost $O(1)$. Swapping elements does not exchange subtrees.

\method{expandExternal(v)} adds two children to a leaf and increases size by two. Set each child's parent to $v$.

The slide's \method{removeExternal(v)} removes leaf $v$ and its parent $p$, replacing $p$ by sibling $s$ (also called removal above an external node):
\par\noindent\begin{minipage}{\linewidth}
\begin{lstlisting}
p = v.parent
s = sibling(v)
if p is root:
    root = s
    s.parent = null
else:
    g = p.parent
    replace g's child reference to p with s
    s.parent = g
size = size - 2
invalidate v and p
\end{lstlisting}
\end{minipage}\par
Assume a full tree and a non-root leaf with a sibling. Updating \method{s.parent} is necessary in the non-root case. This local update costs $O(1)$.

\subsection{Array binary tree}
With zero-based root numbering:
\[
\text{left}(i)=2i+1,\quad \text{right}(i)=2i+2,\quad
\text{parent}(i)=\left\lfloor\frac{i-1}{2}\right\rfloor\ (i>0).
\]
A child exists only if its index is within storage and its slot is occupied. In a complete tree stored in $[0,n)$, child indices must be strictly less than $n$. Node count alone does not identify occupied indices in a sparse tree.

A complete tree uses $O(n)$ consecutive slots. A right-only chain of $n$ nodes ends at index $2^n-2$, requiring array length $2^n-1$. Moving a subtree can require relocating descendants; linked representations avoid these holes and permit local rewiring.

\subsection{General tree and binary encoding}
A general linked node stores an element, parent, and child-reference collection. An ordered general tree can be encoded by binary links:
\begin{itemize}
\item left link: first child in the general tree;
\item right link: next sibling in the general tree.
\end{itemize}
The roots agree, and the binary root has no next sibling. Binary preorder preserves general-tree preorder. Because right now means sibling, binary depth, height, and parent relations must not be read as those of the original hierarchy.

\paragraph{Study checklist}
Classify full/perfect/complete trees; derive leaf/height bounds; produce four traversal orders; justify queue and recursion space; preserve parent references when rewiring; explain why complete trees fit arrays better than sparse trees.

\end{document}
